% GENERATED FILE. Edit canonical lessons, then run npm run generate:books.
% canonical-source-hash: fnv1a64:e1c34b235f01240e
% canonical-lessons: MR-C302-dacakne, MR-C302-thabakne, MR-C302-ksan, MR-C302-tevadhyat, MR-C302-tabadtob

\chapter{In a Split Second}
\label{ch:mr-in-a-split-second}
\hlchaptermodality{302}

\begin{chapteropening}
\textbf{By the end of this chapter:} \emph{I can tell what happened in a sudden moment.}

Complete the last lesson of chapter 302: I can tell what happened in a sudden moment.
\end{chapteropening}

\section[\texorpdfstring{\mr{दचकणे} (dacakṇe)}{दचकणे (dacakṇe)}]{\mr{दचकणे} (dacakṇe) — to be startled}

\label{lesson:MR-C302-dacakne}



\begin{quote}
Before the new one: say the Marathi for to peep, then the Marathi for to observe closely.
\end{quote}

\subsection*{You'll want to know: \mr{दचकणे}}
\textbf{\mr{दचकणे}} — \emph{dacakṇe} — \textquotedblleft{}to be startled\textquotedblright{}.

\begin{quote}
\mr{दार} \mr{अचानक} \mr{वाजलं} \mr{आणि} \mr{मी} \mr{दचकलो}.
\end{quote}

\begin{quote}
\emph{dār acānak vājlaṁ āṇi mī dacaklo.}
\end{quote}

\begin{quote}
The door suddenly banged and I started.
\end{quote}

Past (he …): \mr{तो} \mr{दचकला}, \emph{to dacaklā}.

Intransitive: \mr{आवाजाने} \mr{दचकणे}, to start at a sound; a woman says \mr{मी} \mr{दचकले}. The letter \mr{ऋ} opens \mr{ऋतू}, a season.

\begin{cousinweb}
\textbf{\mr{दचकणे}} (\emph{dacakṇe}) is inherited: CDIAL lists Marathi \emph{dacakṇe}, beside \emph{dackā}, 'fear', under \emph{*dacc}, 'fear' (CDIAL 6126).
\end{cousinweb}

\begin{grammarlens}[title={What you've built}]
One more word for telling what happened.
\end{grammarlens}

\subsection*{Your turn}
\begin{itemize}
\raggedright
  \item \emph{Say it:} \emph{dacakṇe}
  \item \emph{Say it:} \emph{dacakṇe}, once more
  \item \emph{Say it:} \emph{dār acānak vājlaṁ āṇi mī dacaklo.} — The door suddenly banged and I started.
  \item \emph{Recall:} say the Marathi for to peep, then the Marathi for to observe closely, then say \emph{dacakṇe} again
\end{itemize}

\begin{tcolorbox}[breakable,colback=teal!4,colframe=teal!35!black,title={Before you move on}]
What does \mr{दचकणे} mean? (\textquotedblleft{}To be startled\textquotedblright{}.) Say it once more.
\end{tcolorbox}
\section[\texorpdfstring{\mr{थबकणे} (thabakṇe)}{थबकणे (thabakṇe)}]{\mr{थबकणे} (thabakṇe) — to stop short}

\label{lesson:MR-C302-thabakne}



\begin{quote}
Before the new one: say the Marathi for to observe closely, then the Marathi for to be startled.
\end{quote}

\subsection*{You'll want to know: \mr{थबकणे}}
\textbf{\mr{थबकणे}} — \emph{thabakṇe} — \textquotedblleft{}to stop short\textquotedblright{}.

\begin{quote}
\mr{रस्त्यात} \mr{साप} \mr{दिसल्यावर} \mr{ती} \mr{थबकली}.
\end{quote}

\begin{quote}
\emph{rastyāt sāp disalyāvar tī thabaklī.}
\end{quote}

\begin{quote}
When she saw a snake on the road, she stopped short.
\end{quote}

Past (he …): \mr{तो} \mr{थबकला}, \emph{to thabaklā}.

Intransitive: \mr{थबकून} \mr{उभं} \mr{राहणे}, to stand stock-still.

\begin{cousinweb}
\textbf{\mr{थबकणे}} (\emph{thabakṇe}) is a Marathi verb of stopping short; its origin is not settled.
\end{cousinweb}

\begin{grammarlens}[title={What you've built}]
One more word for telling what happened.
\end{grammarlens}

\subsection*{Your turn}
\begin{itemize}
\raggedright
  \item \emph{Say it:} \emph{thabakṇe}
  \item \emph{Say it:} \emph{thabakṇe}, once more
  \item \emph{Say it:} \emph{rastyāt sāp disalyāvar tī thabaklī.} — When she saw a snake on the road, she stopped short.
  \item \emph{Recall:} say the Marathi for to observe closely, then the Marathi for to be startled, then say \emph{thabakṇe} again
\end{itemize}

\begin{tcolorbox}[breakable,colback=teal!4,colframe=teal!35!black,title={Before you move on}]
What does \mr{थबकणे} mean? (\textquotedblleft{}To stop short\textquotedblright{}.) Say it once more.
\end{tcolorbox}
\section[\texorpdfstring{\mr{क्षण} (kṣaṇ)}{क्षण (kṣaṇ)}]{\mr{क्षण} (kṣaṇ) — a moment; an instant}

\label{lesson:MR-C302-ksan}



\begin{quote}
Before the new one: say the Marathi for to be startled, then the Marathi for to stop short.
\end{quote}

\subsection*{You'll want to know: \mr{क्षण}}
\textbf{\mr{क्षण}} — \emph{kṣaṇ} — \textquotedblleft{}a moment; an instant\textquotedblright{}.

\begin{quote}
\mr{एका} \mr{क्षणात} \mr{सगळं} \mr{बदललं}.
\end{quote}

\begin{quote}
\emph{ekā kṣaṇāt sagḷaṁ badallaṁ.}
\end{quote}

\begin{quote}
Everything changed in a moment.
\end{quote}

Masculine.

\mr{क्षणभर}, for a moment; \mr{त्या} \mr{क्षणी}, at that moment.

\begin{cousinweb}
\textbf{\mr{क्षण}} (\emph{kṣaṇ}) is a tatsama, taken from Sanskrit \emph{kṣaṇa}, 'an instant, a moment' (Monier-Williams).
\end{cousinweb}

\begin{grammarlens}[title={What you've built}]
One more word for telling what happened.
\end{grammarlens}

\subsection*{Your turn}
\begin{itemize}
\raggedright
  \item \emph{Say it:} \emph{kṣaṇ}
  \item \emph{Say it:} \emph{kṣaṇ}, once more
  \item \emph{Say it:} \emph{ekā kṣaṇāt sagḷaṁ badallaṁ.} — Everything changed in a moment.
  \item \emph{Recall:} say the Marathi for to be startled, then the Marathi for to stop short, then say \emph{kṣaṇ} again
\end{itemize}

\begin{tcolorbox}[breakable,colback=teal!4,colframe=teal!35!black,title={Before you move on}]
What does \mr{क्षण} mean? (\textquotedblleft{}A moment; an instant\textquotedblright{}.) Say it once more.
\end{tcolorbox}
\section[\texorpdfstring{\mr{तेवढ्यात} (tevaḍhyāt)}{तेवढ्यात (tevaḍhyāt)}]{\mr{तेवढ्यात} (tevaḍhyāt) — at that very moment; meanwhile}

\label{lesson:MR-C302-tevadhyat}



\begin{quote}
Before the new one: say the Marathi for to stop short, then the Marathi for a moment.
\end{quote}

\subsection*{You'll want to know: \mr{तेवढ्यात}}
\textbf{\mr{तेवढ्यात}} — \emph{tevaḍhyāt} — \textquotedblleft{}at that very moment; meanwhile\textquotedblright{}.

\begin{quote}
\mr{मी} \mr{निघणार} \mr{होतो}, \mr{तेवढ्यात} \mr{पाऊस} \mr{सुरू} \mr{झाला}.
\end{quote}

\begin{quote}
\emph{mī nighṇār hoto, tevaḍhyāt pāūs surū jhālā.}
\end{quote}

\begin{quote}
I was about to leave when, at that very moment, it started to rain.
\end{quote}

Built on \mr{तेवढा}, that much, with -\mr{त}: in that much time.

\begin{cousinweb}
\textbf{\mr{तेवढ्यात}} (\emph{tevaḍhyāt}) is built from \emph{tevaḍhā}, 'that much', with the locative ending \emph{-t}: within that much time. The origin of \emph{tevaḍhā} is not settled.
\end{cousinweb}

\begin{grammarlens}[title={What you've built}]
One more word for telling what happened.
\end{grammarlens}

\subsection*{Your turn}
\begin{itemize}
\raggedright
  \item \emph{Say it:} \emph{tevaḍhyāt}
  \item \emph{Say it:} \emph{tevaḍhyāt}, once more
  \item \emph{Say it:} \emph{mī nighṇār hoto, tevaḍhyāt pāūs surū jhālā.} — I was about to leave when, at that very moment, it started to rain.
  \item \emph{Recall:} say the Marathi for to stop short, then the Marathi for a moment, then say \emph{tevaḍhyāt} again
\end{itemize}

\begin{tcolorbox}[breakable,colback=teal!4,colframe=teal!35!black,title={Before you move on}]
What does \mr{तेवढ्यात} mean? (\textquotedblleft{}At that very moment; meanwhile\textquotedblright{}.) Say it once more.
\end{tcolorbox}
\section[\texorpdfstring{\mr{ताबडतोब} (tābaḍtob)}{ताबडतोब (tābaḍtob)}]{\mr{ताबडतोब} (tābaḍtob) — immediately; at once}

\label{lesson:MR-C302-tabadtob}



\begin{quote}
Before the new one: say the Marathi for a moment, then the Marathi for at that very moment.
\end{quote}

\subsection*{You'll want to know: \mr{ताबडतोब}}
\textbf{\mr{ताबडतोब}} — \emph{tābaḍtob} — \textquotedblleft{}immediately; at once\textquotedblright{}.

\begin{quote}
\mr{रुग्णवाहिका} \mr{ताबडतोब} \mr{घटनास्थळी} \mr{पोहोचली}.
\end{quote}

\begin{quote}
\emph{rugṇavāhikā tābaḍtob ghaṭnāsthaḷī pohoclī.}
\end{quote}

\begin{quote}
The ambulance reached the scene at once.
\end{quote}

It keeps one form: \mr{ताबडतोब} \mr{निघणे}, to set off at once.

\begin{cousinweb}
\textbf{\mr{ताबडतोब}} (\emph{tābaḍtob}) is a Marathi word for at once; its origin is not settled.
\end{cousinweb}

\begin{grammarlens}[title={What you've built}]
One more word for telling what happened.
\end{grammarlens}

\subsection*{Your turn}
\begin{itemize}
\raggedright
  \item \emph{Say it:} \emph{tābaḍtob}
  \item \emph{Say it:} \emph{tābaḍtob}, once more
  \item \emph{Say it:} \emph{rugṇavāhikā tābaḍtob ghaṭnāsthaḷī pohoclī.} — The ambulance reached the scene at once.
  \item \emph{Recall:} say the Marathi for a moment, then the Marathi for at that very moment, then say \emph{tābaḍtob} again
\end{itemize}

\begin{tcolorbox}[breakable,colback=teal!4,colframe=teal!35!black,title={Before you move on}]
What does \mr{ताबडतोब} mean? (\textquotedblleft{}Immediately; at once\textquotedblright{}.) Say it once more.
\end{tcolorbox}
